\documentclass[11pt]{article}
% ===========================================================================
%  Shared preamble for the AP Calculus AB exam-preparation worksheets.
%  Input right after \documentclass. Files beginning with "_" are not compiled
%  on their own. Navy and gold, to match the course's lesson pages.
% ===========================================================================
\usepackage[letterpaper,top=2.3cm,bottom=2.2cm,left=1.7cm,right=1.7cm,headheight=28pt,footskip=30pt]{geometry}
\usepackage{amsmath,amssymb}
\usepackage[T1]{fontenc}
\usepackage[utf8]{inputenc}
\usepackage{lmodern}
\usepackage{xcolor}
\usepackage[most]{tcolorbox}
\usepackage{enumitem}
\usepackage{fancyhdr}
\usepackage{array}
\usepackage{tikz}
\usepackage{pgfplots}
\pgfplotsset{compat=1.18}

\definecolor{navy}{HTML}{161E54}
\definecolor{ink}{HTML}{27357E}
\definecolor{gold}{HTML}{B98A00}
\definecolor{goldbg}{HTML}{FFFDF3}
\definecolor{inkbg}{HTML}{F6F7FD}
\definecolor{rule}{HTML}{C9CFF0}
\definecolor{okgreen}{HTML}{1B7A44}
\definecolor{okbg}{HTML}{F1F8F3}

\pagestyle{fancy}
\fancyhf{}
\fancyhead[L]{\footnotesize Name: \underline{\hspace{4cm}}}
\fancyhead[C]{\textbf{\color{navy}Integration Academy}\\[-2pt]{\scriptsize Calculus AB \textperiodcentered{} AP\textsuperscript{\textregistered} Exam Preparation}}
\fancyfoot[L]{\tiny\color{black!55}\parbox[t]{0.78\textwidth}{AP\textsuperscript{\textregistered} is a registered trademark of the College Board, which was not involved in the production of, and does not endorse, this material. \textcopyright{} Integration Academy.}}
\fancyfoot[R]{\scriptsize\thepage}
\renewcommand{\headrulewidth}{1.2pt}
\renewcommand{\headrule}{\hbox to\headwidth{\color{gold}\leaders\hrule height \headrulewidth\hfill}}
\renewcommand{\footrulewidth}{0pt}
\setlength{\parindent}{0pt}
\setlength{\parskip}{4pt}

% Title block: \wstitle{code}{title}{unit line}
\newcommand{\wstitle}[3]{\begin{center}{\small\bfseries\color{gold}WORKSHEET #1}\\[3pt]{\LARGE\bfseries\color{navy}#2}\\[3pt]{\small\color{black!60}#3}\end{center}\vspace{2pt}}
% Section heading
\newcommand{\sech}[1]{\par\vspace{10pt}{\large\bfseries\color{navy}#1}\par\nobreak\vspace{2pt}{\color{rule}\hrule height0.8pt}\vspace{7pt}}
% A level tag printed beside a group of questions
\newcommand{\level}[1]{\par\vspace{6pt}{\footnotesize\bfseries\color{gold}\MakeUppercase{#1}}\par\nobreak\vspace{2pt}}

\newtcolorbox{ideabox}{colframe=ink,colback=inkbg,coltitle=white,colbacktitle=ink,fonttitle=\bfseries\small,title={Key ideas},arc=2pt,boxrule=0.8pt,left=8pt,right=8pt}
\newtcolorbox{example}[1]{colframe=rule,colback=white,coltitle=navy,colbacktitle=inkbg,fonttitle=\bfseries\small,title={#1},arc=2pt,boxrule=0.8pt,borderline west={3pt}{0pt}{ink},left=10pt,right=8pt}
\newtcolorbox{frbox}[1]{colframe=gold,colback=goldbg,coltitle=white,colbacktitle=gold,fonttitle=\bfseries\small,title={#1},arc=2pt,boxrule=0.8pt,breakable,left=8pt,right=8pt}
\newtcolorbox{solbox}{colframe=okgreen,colback=okbg,coltitle=white,colbacktitle=okgreen,fonttitle=\bfseries\small,title={Worked solutions},arc=2pt,boxrule=0.8pt,breakable,left=8pt,right=8pt}
\newtcolorbox{scorebox}[1]{colframe=navy,colback=white,coltitle=white,colbacktitle=navy,fonttitle=\bfseries\small,title={#1},arc=2pt,boxrule=0.8pt,breakable,left=8pt,right=8pt}

\newcommand{\soln}{\par\smallskip\textbf{\color{ink}Solution.}\ }
\newcommand{\result}[1]{\par\smallskip\textbf{#1}}
% blank working space of a given height
\newcommand{\work}[1]{\par\vspace{#1}}
% a point value in the margin of a scoring guide
\newcommand{\pts}[1]{\hfill{\footnotesize\bfseries\color{gold}[#1]}}

% Coordinate grid for graph questions: \begin{apgraph}{xmin}{xmax}{ymin}{ymax} ... \end{apgraph}
\newenvironment{apgraph}[4]{\begin{center}\begin{tikzpicture}\begin{axis}[
  axis lines=middle, xmin=#1, xmax=#2, ymin=#3, ymax=#4,
  grid=both, grid style={black!12}, xtick distance=1, ytick distance=1,
  ticklabel style={font=\scriptsize}, xlabel={$x$}, ylabel={$y$},
  width=10cm, height=7.2cm, samples=80, clip=false,
  every axis plot/.append style={line width=1.1pt, color=ink}]}{\end{axis}\end{tikzpicture}\end{center}}
\newcommand{\closeddot}[2]{\addplot[only marks, mark=*, mark size=2.2pt, color=ink] coordinates {(#1,#2)};}
\newcommand{\opendot}[2]{\addplot[only marks, mark=*, mark size=2.2pt, color=ink, mark options={fill=white}] coordinates {(#1,#2)};}

% Multiple-choice options: four across, or stacked when the options are long.
\newcommand{\opts}[4]{\par\vspace{3pt}\begin{tabular}{@{}p{0.235\linewidth}p{0.235\linewidth}p{0.235\linewidth}p{0.235\linewidth}@{}}(A)\ #1 & (B)\ #2 & (C)\ #3 & (D)\ #4\end{tabular}\par\vspace{9pt}}
\newcommand{\optsv}[4]{\par\vspace{3pt}\begin{tabular}{@{}l@{\ \ }p{0.9\linewidth}@{}}(A) & #1\\(B) & #2\\(C) & #3\\(D) & #4\end{tabular}\par\vspace{9pt}}

\fancyhead[R]{\footnotesize\color{black!60}Worksheet 4.5}
\begin{document}

\wstitle{4.5}{Related Rates: Solving}{Unit 4: Derivatives in Context}

\begin{ideabox}
\begin{itemize}[leftmargin=*,itemsep=2pt,topsep=2pt]
  \item A length that does not change (a ladder, a rope of fixed length) is a \textbf{constant}: put its value in before differentiating.
  \item A length that changes must stay a \textbf{variable} until after differentiating.
  \item For a cone, use similar triangles to write $r$ in terms of $h$ (or the reverse) before differentiating.
  \item For an angle, choose the trigonometric ratio that involves only one changing length if you can.
\end{itemize}
\end{ideabox}

\sech{Worked examples}
\begin{example}{Example 1 \textperiodcentered{} A ladder}
A 17 m ladder leans against a wall. Its foot is pulled away from the wall at 1.5 m/s. How fast is the top moving when the foot is 8 m from the wall?\soln $x^2+y^2=289$, so $x\dfrac{dx}{dt}+y\dfrac{dy}{dt}=0$. When $x=8$, $y=15$: $8(1.5)+15\dfrac{dy}{dt}=0$.\result{The top is sliding down at 0.8 m/s.}
\end{example}

\begin{example}{Example 2 \textperiodcentered{} A cone}
Water is poured at 2 m$^3$/min into an inverted cone of height 6 m and top radius 4 m. How fast is the water level rising when the water is 3 m deep?\soln $r=\dfrac{2h}{3}$, so $V=\dfrac13\pi\left(\dfrac{2h}{3}\right)^2h=\dfrac{4\pi h^3}{27}$ and $\dfrac{dV}{dt}=\dfrac{4\pi h^2}{9}\dfrac{dh}{dt}$. At $h=3$: $2=4\pi\dfrac{dh}{dt}$.\result{$\dfrac{1}{2\pi}$ m/min.}
\end{example}

\begin{example}{Example 3 \textperiodcentered{} A shadow}
A person 1.8 m tall walks at 1.2 m/s away from a lamp post 5.4 m high. How fast is the length of the shadow increasing?\soln With $x$ the distance from the post and $s$ the shadow: $\dfrac{x+s}{5.4}=\dfrac{s}{1.8}$, so $s=\dfrac x2$.\result{$0.6$ m/s.}
\end{example}

\newpage
\sech{Practice}

\level{Core}
\begin{enumerate}[leftmargin=*,itemsep=2pt]
  \item A circular ripple has radius increasing at 2 cm/s. How fast is its area increasing when the radius is 10 cm?\work{2.3cm}
  \item The edge of a cube is increasing at 0.1 cm/s. How fast is its volume increasing when the edge is 5 cm?\work{2.3cm}
  \item A 15 m ladder leans against a wall. Its foot moves away from the wall at 2 m/s. How fast is the top moving when the foot is 9 m from the wall?\work{2.3cm}
  \item A spherical balloon is inflated at $64\pi$ cm$^3$/s. How fast is its radius increasing when $r=2$ cm?\work{2.3cm}
\end{enumerate}

\level{Exam level}
\begin{enumerate}[leftmargin=*,itemsep=2pt,start=5]
  \item Sand is poured at 10 m$^3$/min onto a conical pile whose height always equals its radius. How fast is the height increasing when the pile is 5 m high?\work{2.8cm}
  \item Two cars leave the same point at the same time. One travels east at 60 km/h and the other north at 80 km/h. How fast is the distance between them increasing after 2 hours?\work{2.8cm}
  \item A camera on the ground is 300 m from a launch pad. A rocket rises vertically at 150 m/s. How fast is the camera's angle of elevation changing when the rocket is 400 m high?\work{2.8cm}
  \item Water is pumped at $18\pi$ m$^3$/h into a cylindrical tank of radius 3 m. How fast is the water level rising?\work{2.8cm}
\end{enumerate}

\begin{frbox}{Question 9 \textperiodcentered{} Free-response practice}
A boat is pulled toward a dock by a rope that passes through a ring on the dock 5 m higher than the bow of the boat. The rope is pulled in at 2 m/s. Let $x$ be the horizontal distance from the boat to the dock and $z$ the length of rope between the boat and the ring.
\begin{enumerate}[label=(\alph*),leftmargin=*,itemsep=2pt]
  \item Write an equation relating $x$ and $z$, and differentiate it with respect to time.\work{2.6cm}
  \item How fast is the boat approaching the dock when 13 m of rope are out?\work{2.6cm}
  \item Let $\theta$ be the angle between the rope and the water. How fast is $\theta$ changing at that moment?\work{2.6cm}
\end{enumerate}
\end{frbox}

\newpage
\begin{solbox}
\begin{enumerate}[leftmargin=*,itemsep=4pt]
  \item $\dfrac{dA}{dt}=2\pi r\dfrac{dr}{dt}=40\pi$ cm$^2$/s.
  \item $3(25)(0.1)=7.5$ cm$^3$/s.
  \item $y=12$. $9(2)+12\dfrac{dy}{dt}=0$, so the top slides down at 1.5 m/s.
  \item $64\pi=4\pi(4)\dfrac{dr}{dt}$, so 4 cm/s.
  \item $r=h$, so $V=\dfrac{\pi h^3}{3}$ and $\dfrac{dV}{dt}=\pi h^2\dfrac{dh}{dt}$: $10=25\pi\dfrac{dh}{dt}$, giving $\dfrac{2}{5\pi}$ m/min.
  \item $x=120$, $y=160$, $z=200$. $200\dfrac{dz}{dt}=120(60)+160(80)$, so 100 km/h.
  \item $\tan\theta=\dfrac{y}{300}$, so $\sec^2\theta\,\dfrac{d\theta}{dt}=\dfrac{150}{300}$. With $\sec^2\theta=1+\dfrac{16}{9}=\dfrac{25}{9}$: $\dfrac{d\theta}{dt}=0.18$ rad/s.
  \item $V=9\pi h$, so $18\pi=9\pi\dfrac{dh}{dt}$: 2 m/h.
  \item \begin{enumerate}[label=(\alph*),leftmargin=*,itemsep=3pt]
    \item $z^2=x^2+25$, so $z\dfrac{dz}{dt}=x\dfrac{dx}{dt}$.
    \item $x=12$ and $\dfrac{dz}{dt}=-2$: $13(-2)=12\dfrac{dx}{dt}$, so the boat approaches at $\dfrac{13}{6}$ m/s.
    \item $\sin\theta=\dfrac5z$, so $\cos\theta\,\dfrac{d\theta}{dt}=-\dfrac{5}{z^2}\dfrac{dz}{dt}$. With $\cos\theta=\dfrac{12}{13}$: $\dfrac{12}{13}\dfrac{d\theta}{dt}=\dfrac{10}{169}$, so $\dfrac{d\theta}{dt}=\dfrac{5}{78}$ rad/s (increasing).
  \end{enumerate}
\end{enumerate}
\end{solbox}

\end{document}
